\TieudeT{PHONG TỎA VẬT LÍ}
\subsubsection*{PHẦN I. Thí sinh trả lời câu hỏi từ câu 1 đến câu 18. Mỗi câu hỏi thí sinh chỉ chọn một phương án}
\setcounter{ex}{0}
%mac
%\loai{Đối tượng nghiên cứu}
\begin{ex}%book
Đối tượng nghiên cứu của Vật lí là
\choice
{\True các dạng vận động của vật chất và năng lượng}
{các dạng vận động của sinh vật và năng lượng}
{các dạng chuyển động của chất rắn và chất lỏng}
{các dạng chuyển động của các hành tinh và ngôi sao}

\haidong{6}\end{ex}



%\loai{Ứng dụng của Vật lí}

\begin{ex}%book
Vận dụng kiến thức sự nở vì nhiệt của các chất để chế tạo nhiệt kế rượu, nhiệt kế thủy ngân là ứng dụng Vật lí ở lĩnh vực nào sau đây?
\choice
{Nông nghiệp}
{Công nghiệp}
{\True Y tế}
{Thông tin liên lạc}

\haidong{6}\end{ex}

\begin{ex}%book
Trong các lĩnh vực sau:\\
(1) Vũ trụ.\quad
\quad(2) Kĩ thuật.
\quad(3) Công nghệ.
\quad(4) Nông nghiệp. 
\quad(5) Y tế.	\\
Số lĩnh vực mà Vật lí có ảnh hưởng đến là
\choice
{2}
{3}
{4}
{\True 5}

\haidong{6}\end{ex}


%\loai{Quá trình phát triển}
\begin{ex}%book
Quá trình phát triển của Vật lí được chia thành bao nhiêu giai đoạn chính?
\choice
{\True 3}
{4}
{2}
{5}	
\loigiai{Quan sát - Thực nghiệm - Mô hình}
\haidong{6}\end{ex}


%\loai{Cách mạng công nghiệp 1}
\begin{ex}%book
Máy hơi nước do James Watt sáng chế được ra đời trong cuộc cách mạng công nghiệp lần thứ
\choice
{\True nhất}
{hai}
{ba}
{tư}
\haidong{6}\end{ex}

\begin{ex}%book
Đặc trưng của cuộc cách mạng công nghiệp lần thứ nhất là
\choice
{\True thay thế sức lực cơ bắp bằng máy móc}
{sử dụng các thiết bị điện trong mọi lĩnh vực của đời sống}
{tự động hóa các quá trình sản xuất}
{sử dụng trí tuệ nhân tạo, robot và internet toàn cầu}

\haidong{6}\end{ex}

%\loai{Cách mạng công nghiệp 2}

\begin{ex}%book
Cuộc cách mạng công nghiệp lần thứ hai vào cuối thế kỉ XIX với khám phá hiện tượng cảm ứng điện từ. Một trong những đặc trưng cơ bản của cuộc cách mạng công nghiệp này là
\choice
{thay thế sức lực cơ bắp bằng sức lực máy móc}
{sử dụng trí tuệ nhân tạo, robot, internet toàn cầu, công nghệ vật liệu siêu nhỏ}
{\True sự xuất hiện của các thiết bị dùng điện trong lĩnh vực sản xuất và đời sống}
{tự động hóa các quá trình sản xuất tự động}
\haidong{6}\end{ex}
%\loai{Cách mạng công nghiệp 4}

\begin{ex}%book
Thành tựu Vật lí nào sau đây \textbf{không} thuộc cuộc cách mạng khoa học lần thứ tư?
\choice
{Ô tô không người lái}
{Rôbốt}
{\True Động cơ hơi nước}
{Điện thoại thông minh}
\haidong{6}%<MyLT1>
\end{ex}

%\loai{Phân loại phương pháp nghiên cứu}
\begin{ex}%book
Phương pháp nghiên cứu thường sử dụng của Vật lí
\choice
{phương pháp lí thuyết và phương pháp thu thập số liệu}
{\True phương pháp thực nghiệm và phương pháp mô hình}
{phương pháp thực nghiệm và phương pháp quy nạp}
{phương pháp lí thuyết và phương pháp định tính}	
\haidong{6}\end{ex}
%\loai{Phương pháp thực nghiệm}
\begin{ex}%book
Cho các số tương ứng với các bước:
\begin{enumEX}[1.]{2}
\item Hình thành giả thuyết;		
\item Đề xuất vấn đề;		
\item Quan sát, suy luận;		
\item Kiểm tra giả thuyết;		
\item Rút ra kết luận;		
\end{enumEX}

Tiến trình tìm hiểu thế giới tự nhiên dưới góc độ Vật lí theo thứ tự các bước sau đây:
\choice
{$1-2-3-4-5$}
{$2-1-3-4-5$}
{\True $3-2-1-4-5$}
{$2-3-1-4-5$}

\haidong{6}\end{ex}

\begin{ex}%book
Galileo Galilei thả rơi hai vật có khối lượng khác nhau tại tháp nghiêng Pisa (nước Ý) để nghiên cứu chuyển động của chúng. Phương pháp nghiên cứu của Galileo được gọi là phương pháp
\choice
{\True thực nghiệm}
{mô hình}
{toán học}
{lí thuyết}
\haidong{6}\end{ex}

\begin{ex}%book
Các hiện tượng Vật lí nào sau đây \textbf{không} liên quan đến phương pháp thực nghiệm?
\choice
{\True Tính toán quỹ đạo chuyển động của Thiên Vương Tinh dựa vào toán học}
{Ném một quả bóng lên trên cao}
{Thả rơi một vật từ trên cao xuống mặt đất}
{Kiểm tra sự thay đổi nhiệt độ trong quá trình nóng chảy hoặc bay hơi của một chất}	
\haidong{6}\end{ex}


\begin{ex}%book
Galilei làm thí nghiệm tại tháp nghiêng Pisa vào khoảng năm
\choice
{\True $1600$}
{$1687$}
{$1785$}
{$1831$}	
\haidong{6}%<MyLT1>
\end{ex}


%\loai{Phương pháp mô hình}

\begin{ex}%book
Sắp xếp các bước của phương pháp mô hình theo thứ tự đúng:\\ Kết luận $(1)$, kiểm tra sự phù hợp $(2)$, xác định đối tượng $(3)$, xây dựng mô hình $(4)$.
\choice
{$(1)$, $(2)$, $(3)$, $(4)$}
{\True $(3)$, $(4)$, $(2)$, $(1)$}
{$(4)$, $(3)$, $(2)$, $(1)$}
{$(2)$, $(3)$, $(4)$, $(1)$}
\haidong{6}\end{ex}

\begin{ex}%Tailieuchuan
Mục tiêu của vật lí là
\choice
{khám phá sự vận động của máy móc}
{tìm quy luật về sự chuyển động của các hành tinh}
{\True tìm quy luật chi phối sự vận động của vật chất và năng lượng}
{tìm ra cấu tạo của các chất}
\loigiai{

}
\end{ex}
\begin{ex}%book
Các hiện tượng Vật lí nào sau đây \textbf{không} liên quan đến phương pháp mô hình?
\choice
{Để biểu diễn đường truyền của ánh sáng người ta dùng tia sáng}
{\True Ném một quả bóng lên trên cao}
{Tính toán quỹ đạo chuyển động của Thiên Vương Tinh dựa vào toán học}
{Quả địa cầu là mô hình thu nhỏ của Trái Đất}

\haidong{6}\end{ex}

\begin{ex}%Tailieuchuan
Cấp độ vi mô là
\choice
{\True cấp độ dùng để mô phỏng vật chất nhỏ bé}
{cấp độ to, nhỏ tùy thuộc vào quy mô được khảo sát}
{cấp độ dùng để mô phỏng tầm rộng lớn hay rất lớn của vật chất}
{cấp độ dùng để mô phỏng các vật thể xung quanh đời sống con người}
\loigiai{

}
\end{ex}
\begin{ex}%Tailieuchuan
Người được mệnh danh là "cha đẻ" của phương pháp thực nghiệm là
\choice
{Newton}
{Aristotle}
{\True Galilei}
{Einstein}
\loigiai{

}
\end{ex}
\subsubsection*{PHẦN 2. Thí sinh trả lời từ câu 1 đến câu 4. Trong mỗi ý a), b), c), d) ở mỗi câu, thí sinh chọn đúng hoặc sai.}
\setcounter{ex}{0}
\begin{ex}%Tailieuchuan
Sự ra đời của phương pháp thực nghiệm: Quan niệm, phương pháp nghiên cứu của Aristotle và Galilei.
\choiceTF[t]
{Aristotle dựa vào thực nghiệm để đưa ra quan điểm vật nặng rơi nhanh hơn vật nhẹ} 
{\True Từ những quan sát trong đời sống, Galilei dự đoán: Sự rơi nhanh hay chậm không phụ thuộc vật nặng hay nhẹ} 
{Galilei thiết kế và tiến hành thí nghiệm tại tháp nghiêng Pisa: thả rơi cùng lúc hai quả cầu bằng kim loại, quả to nặng gấp 10 lần quả nhỏ. Kết quả thu được hai quả cầu đều rợi nhanh như nhau, cùng chạm đất cùng một lúc. Kết quả thí nghiệm đã chứng tỏ dự đoán của Galilei là sai} 
{\True Galilei đã dùng phương pháp thực nghiệm để kiểm chứng và kết luận dự đoán của mình} 
\loigiai{ 
\begin{itemchoice}
\itemch
\itemch
\itemch
\itemch
\end{itemchoice}
}
\end{ex}
\begin{ex}%Tailieuchuan
Vật lí học không ngừng phát triển qua các thời kỳ lịch sử, từ giai đoạn cổ điển đến hiện đại, đồng thời có mối liên hệ chặt chẽ với các ngành khoa học khác cũng như với tiến trình công nghiệp hóa. Dựa trên hiểu biết về đặc điểm của từng giai đoạn phát triển trong Vật lí và vai trò của Vật lí trong khoa học – công nghệ.
\choiceTF[t]
{\True Đặc trưng của cuộc cách mạng công nghiệp lần thứ 3 là tự động hóa các quá trình sản xuất} 
{Trong giai đoạn Vật lí cổ điển, các nhà vật lí tập trung vào các mô hình lí thuyết tìm hiểu thế giới vi mô và sử dụng thí nghiệm để kiếm chứng} 
{Phương pháp thực nghiệm và phương pháp lí thuyết là hai phương pháp riêng biệt trong các phương pháp nghiên cứu Vật lí} 
{\True Các lĩnh vực liên môn như Vật lí sinh học, Sinh học lượng tử,... là kết quả của mối quan hệ chặt chẽ của Vật lí với mọi ngành khoa học, trong đó Vật lí được coi là cơ sở của khoa học tự nhiên} 
\loigiai{ 
\begin{itemchoice}
\itemch
\itemch
\itemch
\itemch
\end{itemchoice}
}
\end{ex}
\begin{ex}%Tailieuchuan
Vật lí học nghiên cứu các hiện tượng và quy luật tự nhiên từ cấp độ vi mô đến vĩ mô, góp phần giải thích thế giới vật chất xung quanh.
\choiceTF[t]
{Cấp độ vi mô là cấp độ dùng để mô phỏng vật chất nhỏ bé như cái bàn, cây bút, ô tô,..} 
{\True Cấp độ vĩ mô là cấp độ dùng để mô phỏng tầm rộng lớn hay rất lớn của vật chất như Mặt Trăng, Trái Đất, Mặt Trời,..} 
{Trong nhà trường phổ thông, môn vật lí góp phần quan trọng nhất giúp học sinh hình thành tố chất của một người lãnh đạo trong tương lai} 
{\True Mối quan hệ giữa cường độ dòng điện và hiệu điện thế giữa hai đầu dây dẫn thuộc đối tượng nghiên cứu của môn vật lí} 
\loigiai{ 
\begin{itemchoice}
\itemch
\itemch
\itemch
\itemch
\end{itemchoice}
}
\end{ex}
\begin{ex}%Tailieuchuan
Vật lí học đã trải qua một quá trình phát triển lâu dài, từ những quan sát ban đầu trong thời kì cổ đại đến những lí thuyết hiện đại dựa trên thực nghiệm và mô hình toán học. Mỗi giai đoạn phát triển đều phản ánh sự tiến bộ trong tư duy khoa học và phương pháp nghiên cứu. 
\choiceTF[t]
{\True Có ba giai đoạn trong quá trình phát triển Vật lí, gồm: Tiền Vật lí, Vật lí cổ điển và Vật lí hiện đại} 
{\True Các sai lầm của giai đoạn Tiền Vật lí là do các nhà triết học không tiến hành thực nghiệm kiểm chứng mà chỉ dựa vào quan sát và suy luận chủ quan để kết luận} 
{Einstein là người đặt nền móng đầu tiên trong việc sử dụng phương pháp thực nghiệm để tìm hiểu thế giới tự nhiên} 
{Một số kiến thức vật lí có thể được suy luận từ toán học mà không cần phải kiểm chứng bằng thực nghiệm} 
\loigiai{ 
\begin{itemchoice}
\itemch
\itemch
\itemch
\itemch
\end{itemchoice}
}
\end{ex}
\newpage 
\TieudeT{PHONG TỎA VẬT LÍ 10}
\subsubsection*{PHẦN I. Thí sinh trả lời câu hỏi từ câu 1 đến câu 18. Mỗi câu hỏi thí sinh chỉ chọn một phương án}
\setcounter{ex}{0}
\begin{ex}%book
Đối tượng nào sau đây là một trong những đối tượng nghiên cứu của Vật lí?
\choice
{Nghiên cứu sự trao đổi chất trong cơ thể con người}
{Nghiên cứu sự hình thành và phát triển của các tầng lớp trong xã hội}
{Nghiên cứu về triển vọng phát triển của ngành du lịch nước ta trong giai đoạn tới}
{\True Nghiên cứu về chuyển động cơ học}

\haidong{6}\end{ex}

\begin{ex}%book
Lĩnh vực nào sau đây thuộc phạm vi nghiên cứu của Vật lí?
\choice
{\True Cơ học, quang học, thuyết tương đối}
{Điện học, điện từ học, quy luật di truyền}
{Thuyết tương đối, thuyết tiến hoá, âm học}
{Hội họa, âm học, nhiệt học}

\haidong{6}\end{ex}
\begin{ex}%book
Kiến thức về từ trường Trái Đất nào sau đây được dùng để giải thích đặc điểm nào của loài chim di trú?
\choice
{Làm tổ}
{\True Xác định hướng bay}
{Sinh sản}
{Kiếm ăn}

\haidong{6}\end{ex}
\begin{ex}%book
Phát biểu nào sau đây là \textbf{sai} về ảnh hưởng của Vật lí đến một số lĩnh vực trong đời sống và kĩ thuật?
\choice
{\True Vật lí đem lại cho con người những lợi ích tuyệt vời và không gây ra một ảnh hưởng xấu nào}
{Vật lí ảnh hưởng mạnh mẽ và có tác động làm thay đổi mọi lĩnh vực hoạt động của con người}
{Kiến thức Vật lí trong các phân ngành được áp dụng kết hợp để tạo ra kết quả tối ưu}
{Vật lí là cơ sở của khoa học tự nhiên và công nghệ}
\haidong{6}\end{ex}

\begin{ex}%book
Quá trình nào sau đây là quá trình phát triển của Vật lí?
\choice
{Vật lí cổ điển $\rightarrow$ Vật lí trung đại $\rightarrow$ Vật lí hiện đại}
{\True Tiền Vật lí $\rightarrow$ Vật lí cổ điển $\rightarrow$ Vật lí hiện đại}
{Tiền Vật lí $\rightarrow$ Vật lí trung đại $\rightarrow$ Vật lí hiện đại}
{Tiền Vật lí $\rightarrow$ Vật lí cổ điển $\rightarrow$ Vật lí trung đại}

\haidong{6}\end{ex}
\begin{ex}%book
Thành tựu nghiên cứu nào sau đây của Vật lí được coi là có vai trò quan trọng trong việc mở đầu cho cuộc cách mạng công nghiệp lần thứ nhất?
\choice
{Nghiên cứu về lực vạn vật hấp dẫn}
{\True Nghiên cứu về nhiệt}
{Nghiên cứu về cảm ứng điện từ}
{Nghiên cứu về thuyết tương đối}

\haidong{6}\end{ex}	
\begin{ex}%book
Thành tựu nghiên cứu nào sau đây của Vật lí đóng vai trò quan trọng trong việc mở đầu cuộc cách mạng công nghiệp lần thứ 2?
\choice
{Nghiên cứu về lực hấp dẫn}
{\True Nghiên cứu về hiện tượng cảm ứng điện từ}
{Nghiên cứu về vật liệu nano}
{Nghiên cứu về tự động hóa}

\haidong{6}\end{ex}	
\begin{ex}%book
Đặc trưng cơ bản của cuộc cách mạng công nghiệp lần thứ tư là
\choice
{thay thế sức lực cơ bắp bằng sức lực máy móc}
{\True sử dụng trí tuệ nhân tạo, robot, internet toàn cầu, công nghệ vật liệu nano}
{sự xuất hiện các thiết bị dùng điện trong mọi lĩnh vực sản xuất và đời sống con người}
{tự động hóa các quá trình sản xuất}	
\haidong{6}\end{ex}	
\begin{ex}%book
Ví dụ nào sau đây minh họa cho phương pháp thực nghiệm khi nghiên cứu Vật lí là \textbf{sai}?
\choice
{Galileo thả rơi hai vật có khối lượng khác nhau (cùng hình dạng) từ đỉnh tháp nghiêng Pisa và thấy hai vật rơi chạm đất cùng lúc}
{Archimedes ngâm mình trong bồn nước rồi dựa vào hiện tượng nước trong bồn tắm tràn ra ngoài để tìm ra giải đáp cho việc chiếc vương miện của nhà vua có được làm hoàn toàn từ vàng hay không}
{Để kiểm chứng giả thuyết của Thompson về mô hình cấu tạo nguyên tử, Rutherford đã sử dụng tia alpha gồm các hạt mang điện dương bắn vào các nguyên tử kim loại vàng. Kết quả của thí nghiệm đã bác bỏ giả thuyết của Thompson, đồng thời đã giúp khám phá ra hạt nhân nguyên tử}
{\True Công trình dự đoán sự tồn tại của Hải Vương Tinh trong hệ Mặt Trời vào thế kỉ XIX}
\loigiai{
Sao Hải Vương là hành tinh đầu tiên được tìm thấy bằng tính toán lý thuyết. Dựa vào sự nhiễu loạn bất thường của quỹ đạo Sao Thiên Vương, nhà thiên văn Alexis Bouvard đã kết luận rằng quỹ đạo của nó bị nhiễu loạn do tương tác hấp dẫn với một hành tinh nào đó.
}	
\haidong{6}\end{ex}	
\begin{ex}%book
\immini[thm]{Hình dưới là sơ đồ của phương pháp thực nghiệm nhưng chưa đầy đủ.
Nội dung còn thiếu ở bước số 4 là
\choice
{nghiên cứu tổng quan các dự đoán}
{xây dựng mô hình toán học của dự đoán}
{\True thí nghiệm kiểm tra dự đoán}
{suy luận loại trừ các dự đoán}	}{
\begin{tikzpicture}[draw=Mapcolor, very thick,>=stealth',scale=0.6,transform shape]
\node[rectangle, draw, rounded corners, text width=3.5cm,align = center, minimum height=1.5cm] (a1) at (0,0){1. Xác định vấn đề cần nghiên cứu};
\node[rectangle, draw, rounded corners, text width=3.5cm,align = center, minimum height=1.5cm] (a2) at (3,-2.5){2. Quan sát, thu thập thông tin};
\node[rectangle, draw, rounded corners, text width=3.5cm,align = center, minimum height=1.5cm] (a5) at (-3,-2.5){5. Kết luận};
\node[rectangle, draw, rounded corners, text width=3.5cm,align = center, minimum height=1.5cm] (a3) at (3,-5){3. Đưa ra dự đoán};
\node[rectangle, draw, rounded corners, text width=3.5cm,align = center, minimum height=1.5cm] (a4) at (-3,-5){4. \dotfill};
\draw[->] (a1.east)to[bend left=45](a2.north);
\draw[->] (a5.north)to[bend left=45](a1.west);
\draw[->] (a3.south)to[bend left=45](a4.south);
\draw[->] (a4.north east)to[bend left=45](a3.north west);
\draw[->] (a2.south)to (a3.north);
\draw[->] (a4.north)to (a5.south);
\end{tikzpicture}
}
\haidong{6}\end{ex}	
\begin{ex}%book
Trong phương pháp thực nghiệm thì bước đầu tiên là
\choice
{quan sát, thu thập thông tin}
{thí nghiệm kiểm tra}
{\True xác định vấn đề nghiên cứu}
{đưa ra dự đoán}

\haidong{6}\end{ex}
\begin{ex}%book
Phương pháp mô hình là một trong hai phương pháp nghiên cứu Vật lí. Các loại phương pháp mô hình nào sau đây là thường dùng ở trường Trung học phổ thông?
\choice
{\True Mô hình vật chất, mô hình lí thuyết, mô hình toán học}
{Mô hình vật chất, mô hình thực nghiệm, mô hình toán học}
{Mô hình trực quan, mô hình lí thuyết, mô hình toán học}
{Mô hình trực quan, mô hình thực nghiệm, mô hình toán học}

\haidong{6}\end{ex}
\begin{ex}%Tailieuchuan
Yếu tố nào sau đây là quan trọng nhất dẫn tới việc Aristotle mắc sai lầm khi xác định nguyên nhân làm cho các vật rơi nhanh chậm khác nhau?
\choice
{Không có nhà khoa học nào giúp đỡ ông}
{\True Ông không làm thí nghiệm để kiểm tra quan điểm của mình}
{Khoa học chưa phát triển}
{Ông quá tự tin vào suy luận của mình}
\loigiai{

}
\end{ex}
\begin{ex}%Tailieuchuan
Kiến thức vật lí \textbf{không} giúp chúng ta giải thích các hiện tượng tự nhiên nào sau đây?
\choice
{\True Sự hình thành và phát triển của ý thức con người}
{Hiện tượng nhật thực}
{Sự giãn nở vì nhiệt của vật rắn}
{Cầu vồng xuất hiện sau cơn mưa }
\loigiai{

}
\end{ex}

\begin{ex}%Tailieuchuan
Thành tựu nghiên cứu nào sau đây của Vật lí được coi là có vai trò quan trọng trong việc mở đầu cho cuộc cách mạng công nghệ lần thứ hai?
\choice
{Nghiên cứu về lực vạn vật hấp dẫn}
{Nghiên cứu về nhiệt động lực học}
{\True Nghiên cứu về cảm ứng điện từ}
{Nghiên cứu về thuyết tương đối}
\loigiai{

}
\end{ex}
\begin{ex}%Tailieuchuan
Hiện tượng vật lí nào sau đây liên quan đến phương pháp lí thuyết?
\choice
{\True Ô tô khi chạy đường dài có thể xem ô tô như là một chất điểm}
{Thả rơi một vật từ trên cao xuống mặt đất}
{Kiểm tra sự thay đổi nhiệt độ trong quá trình nóng chảy hoặc bay hơi của một chất}
{Ném một quả bóng lên trên cao}
\loigiai{

}
\end{ex}
\begin{ex}%Tailieuchuan
Các động cơ điện mà con người đang sử dụng như xe đạp điện, xe máy điện,... là kết quả nghiên cứu thuộc chủ yếu những phân ngành Vật lí nào?
\choice
{Quang học và cơ học}
{\True Cơ học và điện học}
{Nhiệt học và điện học}
{Quang học và điện học}
\loigiai{

}
\end{ex}

\begin{ex}%Tailieuchuan
Vai trò nào sau đây \textbf{không} phải là của khoa học tự nhiên trong đời sống?
\choice
{Mở rộng sản xuất, phát triển kinh tế}
{Bảo vệ môi trường, ứng phó với biến đổi khí hậu}
{Bảo vệ sức khỏe và cuộc sống của con người}
{\True Định hướng tư tưởng, phát triển hệ thống chính trị}

\loigiai{

}
\end{ex}









\subsubsection*{PHẦN 2. Thí sinh trả lời từ câu 1 đến câu 4. Trong mỗi ý a), b), c), d) ở mỗi câu, thí sinh chọn đúng hoặc sai.}
\setcounter{ex}{0}
\begin{ex}%Tailieuchuan
Lịch sử loài người đã và đang trải qua bốn cuộc cách mạng công nghiệp trên cơ sở vận dụng các thành tựu nghiên cứu của Vật lí.
\choiceTF[t]
{Đặc trưng của cuộc cách mạng công nghiệp lần thứ nhất là thay thế máy móc bằng sức lực cơ bắp} 
{Con người đang trải qua giai đoạn cuối của cách mạng công nghiệp lần thứ 3, chuẩn bị chuyển sang cuộc cách mạng công nghiệp lần thứ tư} 
{\True Kết quả nghiên cứu về cảm ứng điện từ là cơ sở cho sự hình thành và phát triển của cách mạng công nghiệp lần thứ hai} 
{Việc tự động hóa sản xuất trong cuộc cách mạng công nghiệp lần thứ ba là kết quả của việc vận dụng những thành tựu của Vật lí về Nhiệt động lực học} 
\loigiai{ 
\begin{itemchoice}
\itemch
\itemch
\itemch
\itemch
\end{itemchoice}
}
\end{ex}
\begin{ex}%Tailieuchuan
Các phương pháp nghiên cứu trong Vật lí học bao gồm thực nghiệm, lí thuyết và mô hình. Những phương pháp này phối hợp chặt chẽ, giúp khám phá, kiểm chứng và khái quát các hiện tượng vật lí.
\choiceTF[t]
{\True Phương pháp mô hình giúp đơn giản hoá hiện tượng thực tế để dễ phân tích}  
{\True Phương pháp lí thuyết và thực nghiệm có tính bổ trợ cho nhau} 
{Thí nghiệm về sự xuất hiện dòng điện cảm ứng khi đưa nam châm lại gần/ra xa cuộn dây đồng là một ví dụ của phương pháp lí thuyết} 
{Sử dụng đường sức từ khi khảo sát từ trường của nam châm hay dòng điện là một ví dụ của phương pháp lí thuyết} 
\loigiai{ 
\begin{itemchoice}
\itemch
\itemch
\itemch
\itemch
\end{itemchoice}
}
\end{ex}
\begin{ex}%Tailieuchuan
Môn Vật lí trong nhà trường phổ thông nghiên cứu các hiện tượng, quá trình tự nhiên ở nhiều cấp độ khác nhau của vật chất. Đồng thời, môn học này cũng góp phần rèn luyện năng lực tư duy và hiểu biết khoa học.
\choiceTF[t]
{\True Sự phát triển của thực vật không phải đối tượng nghiên cứu môn vật lí} 
{Mục tiêu học tập môn vật lí giúp học sinh hình thành, phát triển năng lực Toán học} 
{Cấp độ vĩ mô dùng để mô phỏng vật chất nhỏ bé} 
{Sự phân hủy các chất là một lĩnh vực nghiên cứu của Vật lí} 
\loigiai{ 
\begin{itemchoice}
\itemch
\itemch
\itemch
\itemch
\end{itemchoice}
}
\end{ex}
\begin{ex}%Tailieuchuan
Quá trình phát triển của Vật lí hiện nay được chia thành ba giai đoạn lớn: Tiền Vật lí, Vật lí cổ điển và Vật lí hiện đại.
\choiceTF[t]
{Vật lí đang trải qua giai đoạn Vật lí cổ điển} 
{Faraday tìm ra hiện tượng cảm ứng điện từ là thành tựu của Vật lí hiện đại} 
{\True Galilei làm thí nghiệm tại tháp nghiêng Pisa năm 1600 là bước đầu của phương pháp thực nghiệm} 
{\True Joule tìm ra các định luật Nhiệt động lực học là thành tựu của Vật lí cổ điển} 
\loigiai{ 
\begin{itemchoice}
\itemch
\itemch
\itemch
\itemch
\end{itemchoice}
}
\end{ex}
